\scanpage{394}\mspage{133}

donc les foncteurs

\[
\unsure{\lambda|f}:Y\longrightarrow\mathrm{Ch}_\infty(X,a,a)\overset{\mathrm{def}}{=}\Omega(X,a)
\]

correspondent à \uline{certains} \deleted{chemins} \inserted{``lacets''} de $\Omega(X,a)$ (mais lacets à origine non précisée), ils ont un type d'homotopie \inserted{\deleted{\unsure{au sens}} topologique} (``nombre de tours'' autour de la sphère $X$) donné par un entier $n\in\mathbb Z$ \inserted{(le ``degré'' du lacet $\lambda$)}. L'intuition évidente, c'est que l'ensemble des degrés \inserted{$\delta$} qu'on peut obtenir ainsi, avec une catégorie \uline{finie} donnée $Y$ (qui représente \deleted{le \inserted{esp} cercle \unsure{topo}} le cercle $S^1\in(\mathrm{Hot})$), est fini -- \inserted{\unsure{$\lvert\delta\rvert$|!!}} majoré \deleted{\unsure{d'homotopie} $\mathrm{Cat}$ \deleted{\unsure{donc des}}} sûrement par le nombre des \deleted{\unsure{ses|ces}} flèches (donc \unsure{au moins|au mieux} par $4$ dans le cas envisagé, chaque flèche \inserted{de $Y$} faisant faire \uline{au plus} un tour autour de \deleted{\unsure{S}} du pseudo-cercle $\Omega(S^2,a)$).

% slanted
\margin{probablement on \unsure{ne|en} trouve tout au plus que les degrés $\delta$ entre $-2$ et $+2$.}

D'autre part, quand on représente $S^1$ par des catégories $Y_n^{\deleted{\unsure{*}}}$ \deleted{\unsure{multiplicité}} arbitrairement grandes (disons déduites des \deleted{\unsure{\Omega}} intervalles $\unsure{\mathscr L|\mathcal L}_n=(\Delta_+\Delta_-)^n$ en identifiant \deleted{l'origine} \inserted{la} source et le but), on doit trouver des degrés arbitrairement élevés. Or toute $Y_n$ s'envoie dans $Y$ $(=Y_2)$ par une \deleted{\unsure{ff}} \inserted{flèche}

\[
Y_n=Y'\unsure{-|\longrightarrow}Y=Y_2
\]

% archive sheet number: 393
